% Generated draft: requires mathematical and source review.
\documentclass[11pt]{article}
\usepackage[margin=25mm]{geometry}
\usepackage{fontspec}
\setmainfont{Latin Modern Roman}
\usepackage{amsmath,amssymb,mathtools,unicode-math}
\setmathfont{Latin Modern Math}
\usepackage{enumitem,needspace,xurl,hyperref}
\setlength{\parindent}{0pt}
\setlength{\parskip}{5pt}
\setlength{\emergencystretch}{3em}
\newcommand{\N}{\mathbb{N}}
\newcommand{\Z}{\mathbb{Z}}
\newcommand{\Q}{\mathbb{Q}}
\newcommand{\R}{\mathbb{R}}
\newcommand{\C}{\mathbb{C}}
\newcommand{\F}{\mathbb{F}}
\newcommand{\A}{\mathbb{A}}
\newcommand{\PP}{\mathbb{P}}
\newcommand{\E}{\mathbb{E}}
\newcommand{\eps}{\varepsilon}
\newcommand{\dd}{\,\mathrm{d}}
\begin{document}
\section*{UnsolvedMath ID: 30002132}
\textbf{Source problem code:} catalogue entry
\subsection{Definitions and mathematical statement}
For a smooth manifold \(M\), write
\[
\operatorname{crit}(M)=\min_{f\in C^\infty(M,\R)}\#\{x\in M:\dd f_x=0\}.
\]
Here all smooth functions are allowed; the supplied statement does not restrict to Morse functions. The symbol \(\#\) counts distinct points, allowing infinite cardinalities.

To define translated points, suppose a contact form \(\alpha\) defining \(\xi\) has been fixed: on a \((2n+1)\)-dimensional manifold, \(\xi=\ker\alpha\) and \(\alpha\wedge(\dd\alpha)^n\) is nowhere zero. Its Reeb vector field \(R_\alpha\) is determined by
\[
\alpha(R_\alpha)=1,\qquad \dd\alpha(R_\alpha,\cdot)=0.
\]
Write \(\theta_\alpha^s(x)\) for its flow wherever defined. For a contactomorphism \(\phi\) preserving the coorientation, write \(\phi^*\alpha=e^{g_\phi}\alpha\). Its set of translated points relative to \(\alpha\) is
\[
\operatorname{Trans}_\alpha(\phi)=\{x\in M:\exists s\in\R\text{ with }\theta_\alpha^s(x)\text{ defined},\ \phi(x)=\theta_\alpha^s(x),\ g_\phi(x)=0\}.
\]
Different choices of the time \(s\) do not count as different translated points at the same \(x\).

The supplied question is whether every sufficiently \(C^0\)-small contact isotopy \(\{\phi_t\}_{t\in I}\), starting at the identity, satisfies
\[
\forall t\in I,\qquad
\#\operatorname{Trans}_\alpha(\phi_t)\geq\operatorname{crit}(M).
\]
Here \(I\) denotes the isotopy's parameter interval, which the record does not specify. This is a conditional formulation requiring clarification: the source record specifies only \((M,\xi)\), whereas the translated-point definition above requires a chosen contact form. It also supplies no precise neighborhood or uniformity condition for ``\(C^0\)-small,'' and does not state whether \(M\) is closed, connected, coorientable, or allowed to have boundary. No compactness, boundary, support, or nondegeneracy hypotheses have been added. These omissions prevent an unqualified precise conjecture from being extracted.

The supplied dated assessment labels the problem partially solved, reports lower bounds in small and structured cases, and reports counterexamples on the standard sphere to an unrestricted general conjecture. It does not specify the hypotheses of those results or establish which version of the displayed small-isotopy question remains unresolved.
\subsection{Short English statement}
The question asks whether every map in a sufficiently \(C^0\)-small contact isotopy has at least as many distinct translated points as the smallest number of critical points of any smooth real-valued function on the manifold. A contact form, the meaning of smallness, and the intended manifold hypotheses must be clarified. The supplied partially solved label and literature assessment require independent review.
\subsection{Sources}
\begin{itemize}
\item Supplied local UnsolvedMath registry record, permanent ID \(30002132\), title \emph{Translated-Point Bounds for Small Contact Isotopies}, source problem code \texttt{catalogue entry}, in the Oberwolfach Reports Open Problems collection. The proposed year is recorded as 2012; no proposer is identified.
\item Supplied report citation: \emph{Geometric Group Theory, Hyperbolic Dynamics and Symplectic Geometry} (2013), Oberwolfach Reports, DOI 10.4171/owr/2012/35, pp.~2139--2203. \url{https://doi.org/10.4171/owr/2012/35}. This is also the supplied literature reference labeled ``Oberwolfach Report 35.'' The citation's 2013 year and the record's 2012 proposed year are retained without resolving their relationship.
\item Supplied literature reference: Sandon, \emph{Morse estimates for translated points}. \url{https://arxiv.org/abs/1110.0691}.
\item Supplied literature reference: \emph{Contactomorphisms of the sphere without translated points}. \url{https://arxiv.org/abs/2210.11002}.
\item Literature assessment dated exactly \texttt{2026-08-21}: the record labels the problem partially solved, reports Morse-type lower bounds for small and several structured contactomorphisms, and reports sphere counterexamples to the unrestricted general conjecture. This is supplied dated curation, not an independently verified literature assessment; the background explicitly requests independent verification of time-sensitive status and cited claims.
\item The record labels its statement \texttt{corrected\_verified} and says its mathematical content was checked against the official report, with notation, context, and question form normalized. This is a supplied source-verification claim, not verification performed for this draft.
\end{itemize}
\end{document}
